\begin{solapartat}
\punts{0,2} $E = k\,\dfrac{Q}{r^2} \Rightarrow Q = \dfrac{r^2 E}{k}$

% DUBTE: la pauta fa servir R_Terra = 6,40×10^6 m, però l'enunciat dona 6,37×10^6 m.
\punts{0,2} $Q = \dfrac{R^2_\mathrm{Terra}E}{k} = \dfrac{\left(6,40\times 10^{6}\right)^2\cdot 150}{8,99\times 10^{9}} = 6,83\times 10^{5}\ \mathrm{C}$ \punts{0,2}

\punts{0,4} Al estar $\vec{E}$ dirigit cap al centre de la Terra, la càrrega ha de ser negativa.
\end{solapartat}

\begin{solapartat}
\punts{0,5} La càrrega de l'electró és negativa, això fa que el camp elèctric faci una força vertical, contrària al pes.
\[ qE = mg \Rightarrow q = \frac{mg}{E} \]

\punts{0,1} $m = \rho V = \rho\,\frac{4}{3}\pi r^3 = 1000\,\frac{4}{3}\pi\left(18\times 10^{-6}\right)^3 = 2,443\times 10^{-11}\ \mathrm{kg}$

\punts{0,1} $\mathit{Pes} = mg = 2,44\times 10^{-11}\times 9,81 = 2,396\times 10^{-10}\ \mathrm{N}$

% DUBTE: a l'original hi diu «2,40×10^{10}» (errata per 2,40×10^{-10}).
\punts{0,1} $q = \dfrac{\mathrm{Pes}}{E} = \dfrac{2,40\times 10^{10}}{150} = 1,60\times 10^{-12}\ \mathrm{C}$ \ (càrrega negativa)

\punts{0,2} $n = \dfrac{-1,60\times 10^{-12}\ \mathrm{C}}{-1,60\times 10^{-19}\ \mathrm{C}} = 10^{7}$ electrons
\end{solapartat}
